\chapter{Concept Selection}
% this is some introduction about the topic
\section{Pre-class Activities}
\subsection{Activity 1: Generating selection Criteria}
\textbf{Objective:}

The objective of this activity is to walk you through in establishing a set of metrics that can be used to narrow down the number of concepts while improving the overall quality. At the end of this activity, you should be able to generate criteria that can accurately differentiate between different concepts.


    
    \begin{enumerate}
        \item Review the chosen product:
        \begin{enumerate}
        \item By this time, it is expected that you have already familiarized yourself with the chosen product and its specifications.
        \end{enumerate}
        \item Take note of any specific customer needs which are relevant to the product.
        \begin{enumerate}
            \item Refer to the customer needs (Session 5 activity) that have been identified for the chosen product.
        \end{enumerate}
        \item Take note of any specific enterprise needs which are relevant to the product.
        \begin{enumerate}
            \item Refer to the enterprise needs that have been identified for the chosen product, if available.
        \end{enumerate}
        \item Draft the selection criteria.
        \begin{enumerate}
            \item Choose at least one dimension that will serve as selection criteria for the concepts.
            \item The selection criteria should effectively distinguish among the concepts.
            \item Expressed the criteria at a high level of abstraction.
            \item Do not include unimportant criteria in the screening matrix to avoid difficulty in interpreting results.
        \end{enumerate}
    \end{enumerate}


\subsection{Activity 2: Creating a Concept-Screening Matrix} % if available
\textbf{Objective:}
The objective of this activity to walk you through the process of setting up a concept-screening matrix which is critical when evaluating and comparing different concepts based on specific dimensions.

Lets start:
\begin{enumerate}
    \item Gather the identified concepts:
    \begin{enumerate}
        \item Referring to the concepts that were identified during the Session 7 activity.
    \end{enumerate}
    \item Determine the dimensions for the concept-screening matrix:
    \begin{enumerate}
        \item Refer to criteria selected on activity 1.
    \end{enumerate}
    \item Set up the concept-screening matrix:
    \begin{enumerate}
        \item Take a blank A4 paper and orient it horizontally (landscape).
        \item Draw a table with
        \begin{itemize}
            \item columns for the identified concepts and
            \item rows for the selected dimensions.
        \end{itemize}
        \item Label the
        \begin{itemize}
            \item first column as "Concepts" and write the names or labels of the identified concepts in the subsequent rows.
            \item the first row with the selected dimensions.
        \end{itemize}
    \end{enumerate}
\end{enumerate}

\section{In-class Activities}
\subsection{Activity 1: concept-screening matrix}
\textbf{Objective:}
The objective of this activity to walk you through the process of reviewing and consolidating draft matrices in a group setting. By following these instructions, you will work towards reaching a consensus on the concept-screening matrix and proceed to finalize the selection matrix.

\textbf{Procedure:}

\begin{enumerate}
    \item Gather the draft matrices prepared by each team member.
    \item Assess the feasibility of merging the draft matrices or selecting the most promising elements from each.
    \item Collaborate with the group to reach a consensus on the concept-screening matrix.
    \item Once a consensus is reached, proceed to finalize the selection matrix.
\end{enumerate}

\subsubsection{Activity 2: Deciding on the Need for Concept Scoring} % if available
% \textbf{Instructions: Deciding on the Need for Concept Scoring}

\textbf{Objective:}

The objective of this activity to walk you through the process of evaluating the need for concept scoring and making an informed decision as a group. By following these instructions, you will assess whether concept scoring is necessary for your project and understand the justification for not employing concept scoring if applicable.

If your team thinks concept scoring is necessary, they can find the steps for performing it in Chapter 8 under the section "Concept Scoring."

\textbf{Procedure:}
\begin{enumerate}
    \item As a group, discuss the need for concept scoring.
    \begin{enumerate}
        \item Make sure every team member has a clear understanding of the purpose and benefits of concept scoring.
        \end{enumerate}
    \item Evaluate whether increased resolution is necessary to better distinguish between competing concepts?
    \item Justification for not employing concept scoring (if applicable):
    \begin{enumerate}
        \item Provide a statement to justify the decision should the group decides that concept scoring is not necessary.
        \item The reason for can due to the level of differentiation already achieved, time constraints, available resources, or other valid reasons.
    \end{enumerate}
    \item Process flow for concept scoring (if applicable):
    \begin{enumerate}
        \item Refer to the instructions and guidelines outlined in Chapter 8 of our textbook, to conduct concept scoring.
        \end{enumerate}
\end{enumerate}